\documentclass[12pt,a4paper]{article}
\usepackage[utf8]{inputenc}
\usepackage[T1]{fontenc}
\usepackage[polish]{babel}
\usepackage{hyperref}

\title{Temat: Out-of-Domain Text Classification}
\author{Selega Kamil,\\ Sielska Małgorzata,\\ Szczepaniak Katarzyna,\\ Szeląg Magdalena}
\date{}

\begin{document}
\maketitle

\section*{Opis Tematu}
Projekt dotyczy analizy odporności modeli klasyfikacji tekstu na zmianę domeny danych (out-of-domain text classification). Celem jest sprawdzenie, w jakim stopniu modele uczone na jednym typie tekstów potrafią poprawnie klasyfikować dane pochodzące z innego źródła.

W ramach projektu planowane jest przeprowadzenie zadania klasyfikacji sentymentu (pozytywny / negatywny).
Modele zostaną wytrenowane na jednym zbiorze danych (np. recenzje filmowe z IMDb Dataset), a następnie testowane:
\begin{itemize}
    \item na danych z tej samej domeny (in-domain),
    \item na danych z innej domeny (out-of-domain, np. Amazon Reviews).
\end{itemize}

Pozwoli to ocenić zdolność modeli do generalizacji oraz porównać spadek jakości klasyfikacji między środowiskiem in-domain i out-of-domain.
Wyniki będą analizowane za pomocą metryk takich jak accuracy oraz F1-score, a końcowym efektem projektu będzie porównanie skuteczności modeli oraz ocena, czy bardziej zaawansowane architektury lepiej radzą sobie poza domeną treningową.

\section*{Porównywane modele}
Porównane zostaną modele o różnym poziomie złożoności:
\begin{itemize}
    \item Regresja logistyczna + TF-IDF
    \item LSTM
    \item BERT
    \item DistilBERT
\end{itemize}

\section*{Proponowane zbiory danych}
\begin{itemize}
    \item \textbf{IMDb Dataset} (review, rating, pos/neg) - recenzje filmów \cite{imdb_dataset}.
    \item \textbf{Amazon Reviews Dataset} (review title, review text, rating) - recenzje produktów \cite{amazon_reviews}.
    \item \textbf{Yelp Reviews} (text, stars) - recenzje restauracji i usług \cite{yelp_dataset}.
\end{itemize}

\bibliographystyle{plain}
\bibliography{bibliography}

\end{document}
